%============================================================
%  YLYW书法学习技术论文 — 参考文献
%============================================================

\begin{thebibliography}{99}

% === YLYW先验系列 ===
\bibitem{ma2024ylyw}
马兴录, 张国安, 李金函, 等.
YLYW：一种基于《易经》先验符号知识的联邦式神经符号具身决策框架[R].
青岛科技大学信息科学技术学院, 技术报告, 2025.

\bibitem{ma2024zhiji}
马兴录, 李金函, 张国安, 等.
知几学习：从"事后奖惩"到"事先征兆"的具身学习范式革新[R].
青岛科技大学信息科学技术学院, 技术报告, 2025.

\bibitem{ma2024shufa}
马兴录, 李金函, 等.
YLYW汉语言文字处理新范式：观物取象与乘承比应[R].
青岛科技大学信息科学技术学院, 技术报告, 2025.

% === 易经相关 ===
\bibitem{yijing}
周文王, 孔子, 等.
周易[M].
公元前11世纪--前5世纪.

\bibitem{xici}
孔子, 等.
系辞传[M]//周易.
公元前5世纪.

\bibitem{weihongsen}
魏宏森.
系统论的基本规律[J].
自然辩证法研究, 1995.

% === 具身智能 ===
\bibitem{driess2023palm}
Driess D, Xia F, Sajjadi M S M, et al.
PaLM-E: An embodied multimodal language model[C].
International Conference on Machine Learning, 2023.

\bibitem{brohan2023rt2}
Brohan A, Brown N, Carbajal J, et al.
RT-2: Vision-language-action models transfer web knowledge to robotic control[C].
Conference on Robot Learning, 2023.

\bibitem{openvla2024}
Kim M J, Pertsch K, Karamcheti S, et al.
OpenVLA: An open-source vision-language-action model[J].
arXiv preprint arXiv:2406.09246, 2024.

\bibitem{octo2023}
Team Octo, Ghosh D, Walke H, et al.
Octo: An open-source generalist robot policy[C].
Robotics: Science and Systems, 2024.

\bibitem{black2024pi0}
Black K, Brown N, Driess D, et al.
$\pi_0$: A vision-language-action flow model for general robot control[J].
arXiv preprint arXiv:2410.24164, 2024.

% === 机器人书法 ===
\bibitem{yao2004calligraphy}
Yao F, Shao G, Yi J.
Extracting the trajectory of writing brush in Chinese character calligraphy[J].
Engineering Applications of Artificial Intelligence, 2004.

\bibitem{matsui2014calligraphy}
Matsui A, Katsura S.
Calligraphy-stroke learning support system using projector and motion sensor[J].
IEEE International Conference on Mechatronics, 2014.

\bibitem{wu2023calligraphy}
Wu R, Zhou C, Chao F, et al.
A developmental evolutionary learning framework for robotic Chinese calligraphy[J].
IEEE Transactions on Cognitive and Developmental Systems, 2023.

\bibitem{chao2023calligraphy}
Chao F, Lin G, Zheng L, et al.
A learning framework for robotic calligraphy with style transfer[J].
Complex \& Intelligent Systems, 2023.

% === 机器人焊接 ===
\bibitem{chen2023welding}
Chen S B, Lv N.
Research evolution on intelligentized technologies for arc welding process[J].
Journal of Manufacturing Processes, 2014.

\bibitem{xu2022welding}
Xu Y, Lv N, Fang G, et al.
Welding seam tracking in robotic gas metal arc welding[J].
Journal of Manufacturing Processes, 2022.

\bibitem{wang2023welding}
Wang X, Zhou Q, Chen B, et al.
Deep learning-based welding seam recognition and tracking[J].
Robotics and Computer-Integrated Manufacturing, 2023.

% === 神经符号系统 ===
\bibitem{garcez2022ns}
Garcez A, Lamb L C.
Neurosymbolic AI: The 3rd wave[J].
Artificial Intelligence Review, 2023.

\bibitem{yu2023ns}
Yu D, Yang B, Liu D, et al.
A survey of neurosymbolic visual reasoning[J].
arXiv preprint arXiv:2305.07625, 2023.

% === 深度学习与强化学习 ===
\bibitem{levine2018rl}
Levine S, Pastor P, Krizhevsky A, et al.
Learning hand-eye coordination for robotic grasping with deep learning and large-scale data collection[J].
International Journal of Robotics Research, 2018.

\bibitem{kalashnikov2018qt}
Kalashnikov D, Irpan A, Pastor P, et al.
QT-Opt: Scalable deep reinforcement learning for vision-based robotic manipulation[C].
Conference on Robot Learning, 2018.

% === 对比方法 ===
\bibitem{fuzzy2020}
Zadeh L A.
Fuzzy logic[J].
Computer, 1988.

\bibitem{wang2025vla}
Wang Y, et al.
VLA模型在具身智能中的语义泛化与执行泛化瓶颈综述[J].
中国科学：信息科学, 2025.

% === 仿真环境 ===
\bibitem{todorov2012mujoco}
Todorov E, Erez T, Tassa Y.
MuJoCo: A physics engine for model-based control[C].
IEEE/RSJ International Conference on Intelligent Robots and Systems, 2012.

\bibitem{coumans2021pybullet}
Coumans E, Bai Y.
PyBullet, a Python module for physics simulation for games, robotics and machine learning[J].
http://pybullet.org, 2016--2021.

\end{thebibliography}
